
\subsubsection{2.1 LLM Self-Repair}

Chen et al. (2023) introduced self-debugging, where LLMs explain their code to identify mistakes. This work reported 12\% improvement on MBPP using error message feedback. Jiang et al. (2024) extended this with chain-of-explanation before refinement. Li et al. (2024) applied tree search over the repair space.

\subsubsection{2.2 Execution Trace-Based Repair}

TraceFixer \citep{bouzenia2024repairagent} uses runtime traces with divergence point detection. DynaFix \citep{huang2025dynafix} incorporates variable states, control-flow paths, and call stacks. TraceCoder \citep{Huang2026} applies fine-grained trace analysis with multi-agent architectures.

\subsubsection{2.3 Automated Program Repair}

Traditional automated program repair uses fault localization to guide patch generation \citep{dikici2025apr}. Recent work has integrated these techniques with LLMs (Li et al., 2024; Kong et al., 2025; Bouzenia et al., 2024).

\subsubsection{2.4 Gap Addressed}

This work provides the first controlled comparison of five granularity levels (G0-G4) using the same model, benchmark, and prompt template across all conditions.
